\section{The Investigation Agent}
\label{sec:method}

Fig.~\ref{fig:arch} shows the architecture. The agent is a client of an investigation runtime: it leases
tasks, issues tool calls, and submits. It uses the runtime's published client library unchanged and replaces
only the reasoning loop. We describe the tool interface (\S\ref{sec:catalog}), the investigation method
(\S\ref{sec:prompt}), the local safeguards (\S\ref{sec:safeguards}), policy-compatible rationale fields
(\S\ref{sec:alias}), parallel execution (\S\ref{sec:parallel}) and the model drivers (\S\ref{sec:drivers}).

% Agent architecture (TikZ). Included by sections/method.tex.
\begin{figure*}[t]
\centering
\begin{tikzpicture}[
  font=\footnotesize,
  box/.style={draw, rounded corners=2pt, align=center, text width=#1, inner sep=3pt, fill=white, minimum height=9mm},
  box/.default=30mm,
  arr/.style={-{Latex[length=2mm]}, thick},
  darr/.style={-{Latex[length=2mm]}, thick, dashed},
  lab/.style={font=\scriptsize, fill=white, inner sep=1pt}]

% top row
\node[box=32mm, fill=gray!10] (rt) at (0,0) {Investigation runtime\\{\scriptsize competition or local:}\\{\scriptsize task queue, evidence graph}};
\node[box=32mm] (cat) at (5.6,0) {Tool catalog\\{\scriptsize advertised per task;}\\{\scriptsize normalised, aliased}};
\node[box=34mm, fill=cA!8] (llm) at (11.2,0) {LLM driver\\{\scriptsize Claude (thinking, effort, cache)}\\{\scriptsize or OpenAI-compatible}};

% middle row
\node[box=26mm, fill=gray!10] (pool) at (0,-2.6) {Worker pool\\{\scriptsize learned task limit}};
\node[box=58mm, fill=cA!8] (inv) at (5.6,-2.6) {Investigation core\\{\scriptsize \tool{execute}(tool, args) with local safeguards}};
\node[box=26mm, fill=gray!10] (tr) at (11.6,-4.5) {Traces and metrics\\{\scriptsize per call, per run}};

% safeguards
\node[box=19mm] (c1) at (2.45,-4.5) {Response cache};
\node[box=19mm] (c2) at (4.55,-4.5) {Schema validation};
\node[box=19mm] (c3) at (6.65,-4.5) {Grounding check};
\node[box=19mm] (c4) at (8.75,-4.5) {Budget, deadline};
\foreach \c in {c1,c2,c3,c4} \draw[thick, gray] (\c.north) -- (\c.north |- inv.south);

% flows
\draw[arr] (rt) -- node[lab, above=1pt]{alert, tool catalog} (cat);
\draw[arr] (cat) -- node[lab, above=1pt]{tool schemas} (llm);
\draw[arr] ([xshift=-6mm]llm.south) |- node[lab, pos=0.3, right=1pt]{tool call} ([yshift=2mm]inv.east);
\draw[arr] ([yshift=-2mm]inv.east) -| node[lab, pos=0.75, right=1pt]{result or local feedback} ([xshift=6mm]llm.south);
\draw[darr] ([xshift=-14mm]inv.north) |- node[lab, pos=0.35, left=1pt]{spent calls only} ([yshift=-2mm]rt.east);
\draw[arr] (pool) -- (inv);
\draw[arr] (c4.east) -- (tr.west);
\end{tikzpicture}
\caption{Architecture. The runtime advertises a tool catalog with each task; the agent normalises it into
model tool definitions. Every tool call passes through the investigation core, whose local safeguards answer
repeated, malformed, ungrounded or over-budget calls without spending a runtime call; only the remaining calls
reach the runtime (dashed). A worker pool runs tasks in parallel up to the runtime's learned limit, and every
call and outcome is traced.}
\label{fig:arch}
\end{figure*}


\subsection{Catalog-driven tool interface}
\label{sec:catalog}
Each task observation carries the list of enabled tools with their argument schemas. The agent treats this
catalog as authoritative and builds the model's tool definitions from it on every task, rather than from a
hard-coded list, so changes to the runtime's tools need no code change. Schemas are normalised for portability
across model providers: local \texttt{\$ref} references are inlined, presentation-only keys are dropped,
\texttt{oneOf} alternatives that are distinguished by a constant discriminator are rewritten as the equivalent
\texttt{anyOf}, and the root is forced to be an object. When a catalog entry lacks a schema, the definition
falls back to the reference client's typed models. A JSON Schema validator (draft 2020-12) is compiled from
each normalised schema.

\subsection{Investigation method}
\label{sec:prompt}
The system prompt (Appendix~\ref{app:prompts}) encodes a hypothesis-driven method in three parts.
\emph{Method:} read the alert fully, form two or three competing hypotheses that always include a benign
explanation (administration, software update, backup, monitoring), and choose each call for the hypothesis it
can settle most decisively, usually by following relations outward from the trigger entity.
\emph{Spending discipline:} before every call, ask whether its answer could change the verdict or the artifact
list; never fetch what the alert or an earlier result already shows; request independent lookups together;
stop as soon as the verdict and the artifacts are supported. \emph{Submission guide:} which response kind fits
which entity (compromised host, used account, persistence, dropped file, control channel, observed entity),
and, for benign verdicts, which kinds of property prove legitimacy (valid signature, expected install path,
known management parent, documented service behaviour). The ablation \emph{no method} replaces this prompt
with a minimal one that states only the tool rules, close to the reference client's instructions.

\subsection{Local safeguards}
\label{sec:safeguards}
Every tool call the model issues is executed by Algorithm~\ref{alg:execute}. Its premise is that a runtime call
should be spent only when its result can be new information; everything else is answered locally, at no cost,
with feedback the model can act on.

\begin{algorithm}[t]
\caption{Execute a model-issued tool call}
\label{alg:execute}
\small
\begin{algorithmic}[1]
\Require tool $t$, arguments $x$; cache $\mathcal C$; rejected set $\mathcal R$; evidence corpus $\mathcal E$; calls spent $n$; hard budget $B$; finish-by time $T$
\State $x \gets \textsc{Unalias}(\textsc{NormaliseJSON}(x))$ \Comment{\S\ref{sec:alias}}
\State $k \gets (t, x \setminus \{\text{rationale}\})$
\If{$t\neq$ finish \textbf{and} $k\in\mathcal C$} \Return cached result \Comment{cache: no call spent}\EndIf
\If{$t\neq$ finish \textbf{and} ($n\ge B$ \textbf{or} now $\ge T$)} \Return ``submit now'' \Comment{budget, deadline}\EndIf
\If{$k\notin\mathcal R$}
  \If{$\textsc{Validate}(t,x)$ fails} $\mathcal R\gets\mathcal R\cup\{k\}$; \Return schema errors \EndIf
  \If{$t=$ finish \textbf{and} some cited id $\notin\mathcal E$} $\mathcal R\gets\mathcal R\cup\{k\}$; \Return ungrounded ids \EndIf
\EndIf
\State $y\gets$ \textsc{CallRuntime}$(t,x)$ with backoff on capacity refusals; $n\gets n+1$
\State $\mathcal E\gets\mathcal E\cup\{y\}$; \textbf{if} $y$ is not an error \textbf{then} $\mathcal C[k]\gets y$
\State \Return \textsc{Truncate}$(y)$
\end{algorithmic}
\end{algorithm}

\paragraph{Response cache} Identical calls (same tool, same arguments, ignoring the free-text rationale) are
answered from a cache with an explicit note that no call was spent. LLM agents re-request evidence they have
already seen, especially across long contexts; the cache makes that free.

\paragraph{Schema validation} Arguments are validated locally against the advertised schema. Invalid calls
receive the validator's messages instead of costing a runtime call. A call rejected locally is sent unchanged
if the model repeats it, so the runtime, not the agent, remains the authority on validity. A related
normalisation unwraps object-typed arguments that a model serialised as JSON strings.

\paragraph{Grounding check} Before a submission is sent, every entity and relation identifier it cites is
checked against the corpus of everything the alert and earlier tool results contained. Identifiers that never
appeared are reported back once (``use ids exactly as returned, or drop them''); an unchanged resubmission is
sent. This targets a failure specific to LLM agents: plausible but invented identifiers in an incident
response.

\paragraph{Budgets and deadline} After each batch of results the model sees the number of calls spent. Past a
soft budget it is told it is outside the efficient range; at a hard budget evidence calls are refused locally,
so the agent submits with what it has instead of exhausting the runtime's step limit and failing the task. An
optional deadline stops taking new tasks a few minutes before a cut-off and, shortly before it, allows only
submissions, so no task is left unsubmitted.

\subsection{Policy-compatible rationale fields}
\label{sec:alias}
The runtime requires a free-text field named \texttt{reasoning} on every evidence call and in the submission.
In early runs with a recent frontier model, the first request of every task was declined before any output,
with the refusal category \texttt{reasoning\_extraction}: the combination of a prompt asking the model to
``state its hypotheses'' and a required field called \texttt{reasoning} described as the ``investigation
trace'' read as a request to externalise the model's internal reasoning, which the model's usage policy
declines. The fix keeps the runtime contract intact. The instruction to state hypotheses was removed from
the first message (the method still asks the model to consider them), and the agent presents the field to the
model as
\texttt{purpose} (``one short sentence naming what this call checks'') and the submission's field as
\texttt{summary} (``a short incident summary''), and maps both back to \texttt{reasoning} before validation
and sending. The model's actual reasoning is read from its summarised thinking blocks. After the change, all
tasks of a mock run completed where all had previously been refused. We report this because any agent built
for a fixed runtime contract can meet the same conflict between the contract and a model provider's policies,
and because a field rename is enough to resolve it.

\subsection{Parallel execution with a learned limit}
\label{sec:parallel}
Budget-scored evaluations are often also time-boxed. The agent runs $W$ workers that each lease and investigate
tasks. Before leasing, a worker reserves a slot, so no more leases are attempted than the current limit allows.
If the runtime refuses a lease because of its active-task limit while other workers are busy, the pool lowers
its limit to the number of tasks in flight and the worker waits for a free slot; refusals while nothing else
runs are waited out with a bounded number of retries. Short operation-capacity or rate-limit refusals on tool
calls are retried with backoff; a refused call was not executed, so retrying spends nothing. One task's failure
is contained to that task.

\subsection{Model drivers}
\label{sec:drivers}
The Claude driver uses native tool use~\cite{anthropic_tooluse} with streaming, adaptive thinking (summarised
thinking is logged to the trace), an explicit effort level, automatic prompt caching~\cite{anthropic_promptcaching}
of the growing conversation prefix, and server-side fallback to another model on policy declines. Assistant turns are appended unchanged, so thinking blocks remain
valid for the next turn, and all results of a batch of parallel tool calls are returned in one message. An
OpenAI-compatible driver supports other providers; malformed tool-call JSON is answered locally without
spending a call. Both drivers nudge the model at most twice when it answers without a tool call, and abort the
task on a refusal. Every task produces a trace: each call with its stated purpose, its status
(\emph{ok}, \emph{cached}, \emph{refused locally}, \emph{error}, \emph{submitted}), the result, the summarised
reasoning, the submission, the runtime's score and token usage. Each run records its configuration and
aggregate metrics.
